\documentclass{article}
\usepackage[utf8]{inputenc}
\usepackage{geometry}
\usepackage{enumitem}
\usepackage{hyperref}
\usepackage{xcolor}
\usepackage{graphicx}
\usepackage{array}
\usepackage{multirow}
\usepackage{amsmath, amssymb}
\usepackage{chemfig}
\usepackage{mhchem}
\usepackage{tikz}
\usepackage{setspace}
\usepackage{soul}
\usepackage{mathtools}
\usepackage{booktabs}
\usepackage{chemformula}
\usepackage{siunitx}
\usetikzlibrary{positioning, calc}
\usepackage{longtable}
\geometry{margin=1in}
\setlength{\parskip}{0.5em}
\usepackage{cancel}


\begin{document}

\begin{center}
    \vspace*{-1cm}
    \begin{tabular}{p{12cm} r}
        & \textbf{Name:} \underline{\hspace{5cm}} \\
    \end{tabular}

    \vspace{0.2cm}
    {\LARGE \textbf{Week 15.1 Worksheet}}
\end{center}    

\noindent\begin{tabular}{|p{0.9\linewidth}|}
\hline
\textbf{(1)} A Brönsted-Lowry acid is defined as a substance that: \\[0.2cm]
\begin{minipage}[t]{0.85\linewidth}
\begin{enumerate}[label=\alph*.]
\item Donates a proton
\item Accepts a proton
\item Releases electrons
\item Accepts electrons
\end{enumerate}
\end{minipage}\\[0.2cm]
\hline
\end{tabular}

\noindent\begin{tabular}{|p{0.9\linewidth}|}
\hline
\textbf{(2)} What is the oxidation number of chlorine in NaCl?\\[0.2cm]
\begin{minipage}[t]{0.85\linewidth}
\begin{enumerate}[label=\alph*.]
\item +1
\item -1
\item 0
\item +2
\end{enumerate}
\end{minipage}\\[0.2cm]
\hline
\end{tabular}

\noindent\begin{tabular}{|p{0.9\linewidth}|}
\hline
\textbf{(3)} Which of the following reactions is a redox reaction?\\[0.2cm]
\begin{minipage}[t]{0.85\linewidth}
\begin{enumerate}[label=\alph*.]
\item Zn + CuSO$_4$ $\rightarrow$ ZnSO$_4$ + Cu\\[0.1cm]
\item AgNO$_3$ + NaCl $\rightarrow$ AgCl + NaNO$_3$\\[0.1cm]
\item CaCO$_3$ $\rightarrow$ CaO + CO$_2$\\[0.1cm]
\item H$_2$O + SO$_3$ $\rightarrow$ H$_2$SO$_4$\\[0.1cm]
\end{enumerate}
\end{minipage}\\[0.2cm]
\hline
\end{tabular}

\noindent\begin{tabular}{|p{0.9\linewidth}|}
\hline
\textbf{(4)} In the reaction $\text{H}_2\text{PO}_4^- \rightleftharpoons \text{HPO}_4^{2-} + \text{H}^+$, the conjugate base is: \\[0.2cm]
\begin{minipage}[t]{0.85\linewidth}
\begin{enumerate}[label=\alph*.]
\item H$_2$PO$_4^-$
\item HPO$_4^{2-}$
\item H$^+$
\item None of the above
\end{enumerate}
\end{minipage}\\[0.2cm]
\hline
\end{tabular}
\newpage

\noindent\begin{tabular}{|p{0.9\linewidth}|}
\hline
\textbf{(5)} For the equilibrium $\text{C}(s) + \text{CO}_2(g) \rightleftharpoons 2\text{CO}(g)$, what happens when pressure increases?\\[0.2cm]
\begin{minipage}[t]{0.85\linewidth}
\begin{enumerate}[label=\alph*.]
\item The system shifts left because fewer gas moles are on the CO$_2$ side
\item The system shifts right because more gas moles are produced
\item No effect because solids and gases respond equally
\item Pressure changes only affect liquids, not gases
\end{enumerate}
\end{minipage}\\[0.2cm]
\hline
\end{tabular}


\noindent\begin{tabular}{|p{0.9\linewidth}|}
\hline
\textbf{(6)} What is the oxidation number of sulfur in SO$_2$?\\[0.2cm]
\begin{minipage}[t]{0.85\linewidth}
\begin{enumerate}[label=\alph*.]
\item +2
\item +4
\item +6
\item 0
\end{enumerate}
\end{minipage}\\[0.2cm]
\hline
\end{tabular}

\noindent\begin{tabular}{|p{0.9\linewidth}|}
\hline
\textbf{(7)} In the reaction $\text{HCl} + \text{H}_2\text{O} \rightarrow \text{H}_3\text{O}^+ + \text{Cl}^-$, which is the Brönsted-Lowry acid?\\[0.2cm]
\begin{minipage}[t]{0.85\linewidth}
\begin{enumerate}[label=\alph*.]
\item HCl
\item H$_2$O
\item H$_3$O$^+$
\item Cl$^-$
\end{enumerate}
\end{minipage}\\[0.2cm]
\hline
\end{tabular}

\noindent\begin{tabular}{|p{0.9\linewidth}|}
\hline
\textbf{(8)} For the exothermic reaction $2\text{SO}_2 + \text{O}_2 \rightarrow 2\text{SO}_3$, increasing temperature will: \\[0.2cm]
\begin{minipage}[t]{0.85\linewidth}
\begin{enumerate}[label=\alph*.]
\item Increase SO$_3$ production
\item Shift equilibrium to left
\item Not affect equilibrium
\item Increase both forward and reverse rate equally
\end{enumerate}
\end{minipage}\\[0.2cm]
\hline
\end{tabular}
\newpage

\noindent\begin{tabular}{|p{0.9\linewidth}|}
\hline
\textbf{(9)} Calcium hydroxide reacts with hydrochloric acid to form calcium chloride and water. Which statement best describes the acid-base behavior occurring in this reaction?\\[0.2cm]
Ca(OH)$_2$ + 2HCl $\rightarrow$ CaCl$_2$ + 2H$_2$O\\[0.2cm]
\begin{minipage}[t]{0.85\linewidth}
\begin{enumerate}[label=\alph*.]
\item HCl donates protons to OH$^-$ ions, forming water as Ca$^{2+}$ pairs with Cl$^-$
\item Ca(OH)$_2$ donates protons to HCl, forming hydrogen gas
\item Both reactants accept protons, preventing water formation
\item Ca$^{2+}$ acts as the proton donor while Cl$^-$ acts as the proton acceptor
\end{enumerate}
\end{minipage}\\[0.2cm]
\hline
\end{tabular}

\noindent\begin{tabular}{|p{0.9\linewidth}|}
\hline
\textbf{(10)} For the reaction $\text{R} + \text{S} \rightleftharpoons \text{RS}$, which change could temporarily increase Q?\\[0.2cm]
\begin{minipage}[t]{0.85\linewidth}
\begin{enumerate}[label=\alph*.]
\item Adding RS
\item Adding R
\item Removing S
\item Diluting all species equally
\end{enumerate}
\end{minipage}\\[0.2cm]
\hline
\end{tabular}

\noindent\begin{tabular}{|p{0.9\linewidth}|}
\hline
\textbf{(11)} For the reaction $2\text{A} \rightleftharpoons \text{B}$, which expression is correct?\\[0.2cm]
\begin{minipage}[t]{0.85\linewidth}
\begin{enumerate}[label=\alph*.]
\item $K_{eq}=\frac{[\text{B}]}{[\text{A}]^2}$
\item $K_{eq}=\frac{[\text{A}]^2}{[\text{B}]}$
\item $K_{eq}=\frac{1}{[\text{A}][\text{B}]}$
\item $K_{eq}=[\text{A}][\text{B}]$
\end{enumerate}
\end{minipage}\\[0.2cm]
\hline
\end{tabular}

\noindent\begin{tabular}{|p{0.9\linewidth}|}
\hline
\textbf{(12)} Identify the reaction type:\\[0.2cm]
BaCl$_2$ + Na$_2$SO$_4$ $\rightarrow$ BaSO$_4$ + 2NaCl\\[0.2cm]
\begin{minipage}[t]{0.85\linewidth}
\begin{enumerate}[label=\alph*.]
\item Combination
\item Decomposition
\item Exchange
\item Single replacement
\end{enumerate}
\end{minipage}\\[0.2cm]
\hline
\end{tabular}

\noindent\begin{tabular}{|p{0.9\linewidth}|}
\hline
\textbf{(13)} In a redox reaction, the oxidizing agent is the species that:\\[0.2cm]
\begin{minipage}[t]{0.85\linewidth}
\begin{enumerate}[label=\alph*.]
\item Loses electrons
\item Gains electrons
\item Is always neutral
\item Has the highest atomic mass
\end{enumerate}
\end{minipage}\\[0.2cm]
\hline
\end{tabular}

\noindent\begin{tabular}{|p{0.9\linewidth}|}
\hline
\textbf{(14)} For the endothermic reaction $\text{N}_2\text{O}_4 \rightarrow 2\text{NO}_2$, decreasing temperature favors: \\[0.2cm]
\begin{minipage}[t]{0.85\linewidth}
\begin{enumerate}[label=\alph*.]
\item Formation of NO$_2$
\item Formation of N$_2$O$_4$
\item Equal amounts of each
\item Faster formation of NO$_2$ only
\end{enumerate}
\end{minipage}\\[0.2cm]
\hline
\end{tabular}
\newpage

\noindent\begin{tabular}{|p{0.9\linewidth}|}
\hline
\textbf{(15)} For the equilibrium $ \text{PCl}_5(g) \rightleftharpoons \text{PCl}_3(g) + \text{Cl}_2(g)$, increasing pressure causes the reaction to shift: \\[0.2cm]
\begin{minipage}[t]{0.85\linewidth}
\begin{enumerate}[label=\alph*.]
\item Left, toward PCl$_5$
\item Right, toward PCl$_3$ and Cl$_2$
\item No shift because gases are unaffected by pressure
\item Toward whichever side has the higher total mass
\end{enumerate}
\end{minipage}\\[0.2cm]
\hline
\end{tabular}


\noindent\begin{tabular}{|p{0.9\linewidth}|}
\hline
\textbf{(16)} In the reaction 2Na + Cl$_2$ $\rightarrow$ 2NaCl, which is the oxidizing agent?\\[0.2cm]
\begin{minipage}[t]{0.85\linewidth}
\begin{enumerate}[label=\alph*.]
\item Na
\item Cl$_2$
\item NaCl
\item Both Na and Cl$_2$
\end{enumerate}
\end{minipage}\\[0.2cm]
\hline
\end{tabular}
\newpage

\noindent\begin{tabular}{|p{0.9\linewidth}|}
\hline
\textbf{(17)} For the reaction $\text{A} + \text{C} \rightleftharpoons \text{AC}$ with $K_{eq} = 0.02$, which is true?\\[0.2cm]
\begin{minipage}[t]{0.85\linewidth}
\begin{enumerate}[label=\alph*.]
\item Reactants are strongly favored
\item Products are strongly favored
\item Equal amounts of A and AC form
\item Keq must be recalculated using solids only
\end{enumerate}
\end{minipage}\\[0.2cm]
\hline
\end{tabular}

\noindent\begin{tabular}{|p{0.9\linewidth}|}
\hline
\textbf{(18)} According to the Arrhenius definition, which statement correctly describes an acid and a base?\\[0.2cm]
\begin{minipage}[t]{0.85\linewidth}
\begin{enumerate}[label=\alph*.]
\item An acid produces H$^+$ in water, and a base produces OH$^-$ in water
\item An acid donates a proton, and a base accepts a proton
\item An acid accepts an electron pair, and a base donates an electron pair
\item An acid and base both increase the concentration of ions equally
\end{enumerate}
\end{minipage}\\[0.2cm]
\hline
\end{tabular}


\noindent\begin{tabular}{|p{0.9\linewidth}|}
\hline
\textbf{(19)} What type of reaction occurs when methane burns in oxygen to produce carbon dioxide and water?\\[0.2cm]
CH$_4$ + 2O$_2$ $\rightarrow$ CO$_2$ + 2H$_2$O\\[0.2cm]
\begin{minipage}[t]{0.85\linewidth}
\begin{enumerate}[label=\alph*.]
\item Combination
\item Decomposition
\item Combustion
\item Single displacement
\end{enumerate}
\end{minipage}\\[0.2cm]
\hline
\end{tabular}

\noindent\begin{tabular}{|p{0.9\linewidth}|}
\hline
\textbf{(20)} What type of chemical reaction occurs between sodium phosphate and calcium nitrate to form calcium phosphate and sodium nitrate?\\[0.2cm]
2Na$_3$PO$_4$ + 3Ca(NO$_3$)$_2$ $\rightarrow$ Ca$_3$(PO$_4$)$_2$ + 6NaNO$_3$\\[0.2cm]
\begin{minipage}[t]{0.85\linewidth}
\begin{enumerate}[label=\alph*.]
\item Decomposition
\item Exchange
\item Combination
\item Combustion
\end{enumerate}
\end{minipage}\\[0.2cm]
\hline
\end{tabular}
\newpage

\noindent\begin{tabular}{|p{0.9\linewidth}|}
\hline
\textbf{(21)} In the reaction $\text{CO}(g) + 2\text{H}_2(g) \rightleftharpoons \text{CH}_3\text{OH}(g)$, which change increases methanol production? \\[0.2cm]
\begin{minipage}[t]{0.85\linewidth}
\begin{enumerate}[label=\alph*.]
\item Increasing pressure because the product side has fewer gas moles
\item Decreasing pressure because the reactant side has more gas moles
\item Increasing volume of the container
\item Adding more inert gas at constant volume
\end{enumerate}
\end{minipage}\\[0.2cm]
\hline
\end{tabular}


\noindent\begin{tabular}{|p{0.9\linewidth}|}
\hline
\textbf{(22)} If $Q > K_{eq}$ for the reaction $\text{F} \rightleftharpoons \text{G}$, which direction will the system shift?\\[0.2cm]
\begin{minipage}[t]{0.85\linewidth}
\begin{enumerate}[label=\alph*.]
\item Toward reactants
\item Toward products
\item No shift
\item The reaction stops
\end{enumerate}
\end{minipage}\\[0.2cm]
\hline
\end{tabular}

\noindent\begin{tabular}{|p{0.9\linewidth}|}
\hline
\textbf{(23)} When sodium reacts with chlorine to produce sodium chloride, what type of reaction is this?\\[0.2cm]
2Na + Cl$_2$ $\rightarrow$ 2NaCl\\[0.2cm]
\begin{minipage}[t]{0.85\linewidth}
\begin{enumerate}[label=\alph*.]
\item Decomposition
\item Combustion
\item Exchange
\item Combination
\end{enumerate}
\end{minipage}\\[0.2cm]
\hline
\end{tabular}

\noindent\begin{tabular}{|p{0.9\linewidth}|}
\hline
\textbf{(24)} For the reaction $\text{P} \rightleftharpoons \text{Q}$, which condition must be true at equilibrium?\\[0.2cm]
\begin{minipage}[t]{0.85\linewidth}
\begin{enumerate}[label=\alph*.]
\item Rate forward = rate reverse
\item $[P]$ = [Q]
\item $K_{eq}$ = 0
\item All reactant is converted to product
\end{enumerate}
\end{minipage}\\[0.2cm]
\hline
\end{tabular}
\newpage

\noindent\begin{tabular}{|p{0.9\linewidth}|}
\hline
\textbf{(25)} For the reaction $\text{A} + \text{heat} \rightarrow \text{B}$, which statement is correct? \\[0.2cm]
\begin{minipage}[t]{0.85\linewidth}
\begin{enumerate}[label=\alph*.]
\item Adding heat increases [B]
\item Adding heat decreases [B]
\item Removing heat increases [B]
\item Heat has no effect on equilibrium
\end{enumerate}
\end{minipage}\\[0.2cm]
\hline
\end{tabular}

\noindent\begin{tabular}{|p{0.9\linewidth}|}
\hline
\textbf{(26)} The burning of benzene in air to form carbon dioxide and water is an example of what type of reaction?\\[0.2cm]
2C$_6$H$_6$ + 15O$_2$ $\rightarrow$ 12CO$_2$ + 6H$_2$O\\[0.2cm]
\begin{minipage}[t]{0.85\linewidth}
\begin{enumerate}[label=\alph*.]
\item Oxidation-reduction
\item Combustion
\item Decomposition
\item Neutralization
\end{enumerate}
\end{minipage}\\[0.2cm]
\hline
\end{tabular}

\noindent\begin{tabular}{|p{0.9\linewidth}|}
\hline
\textbf{(27)} Which element is being reduced in the reaction Fe$_2$O$_3$ + 3CO $\rightarrow$ 2Fe + 3CO$_2$?\\[0.2cm]
\begin{minipage}[t]{0.85\linewidth}
\begin{enumerate}[label=\alph*.]
\item Fe
\item O
\item C
\item CO$_2$
\end{enumerate}
\end{minipage}\\[0.2cm]
\hline
\end{tabular}

\noindent\begin{tabular}{|p{0.9\linewidth}|}
\hline
\textbf{(28)} What happens to an atom's oxidation number when it is oxidized?\\[0.2cm]
\begin{minipage}[t]{0.85\linewidth}
\begin{enumerate}[label=\alph*.]
\item It increases
\item It decreases
\item It stays the same
\item It becomes zero
\end{enumerate}
\end{minipage}\\[0.2cm]
\hline
\end{tabular}
\newpage

\noindent\begin{tabular}{|p{0.9\linewidth}|}
\hline
\textbf{(29)} When ammonium chloride reacts with silver nitrate to form silver chloride and ammonium nitrate, the reaction is classified as:\\[0.2cm]
NH$_4$Cl + AgNO$_3$ $\rightarrow$ AgCl + NH$_4$NO$_3$\\[0.2cm]
\begin{minipage}[t]{0.85\linewidth}
\begin{enumerate}[label=\alph*.]
\item Combination
\item Exchange
\item Decomposition
\item Combustion
\end{enumerate}
\end{minipage}\\[0.2cm]
\hline
\end{tabular}

\noindent\begin{tabular}{|p{0.9\linewidth}|}
\hline
\textbf{(30)} What type of reaction occurs when potassium iodide reacts with lead(II) nitrate to form lead(II) iodide and potassium nitrate?\\[0.2cm]
2KI + Pb(NO$_3$)$_2$ $\rightarrow$ 2KNO$_3$ + PbI$_2$\\[0.2cm]
\begin{minipage}[t]{0.85\linewidth}
\begin{enumerate}[label=\alph*.]
\item Exchange
\item Combination
\item Decomposition
\item Combustion
\end{enumerate}
\end{minipage}\\[0.2cm]
\hline
\end{tabular}

\noindent\begin{tabular}{|p{0.9\linewidth}|}
\hline
\textbf{(31)} When hydrochloric acid reacts with sodium hydroxide to form sodium chloride and water, which statement best describes the acid-base process taking place?\\[0.2cm]
HCl + NaOH $\rightarrow$ NaCl + H$_2$O\\[0.2cm]
\begin{minipage}[t]{0.85\linewidth}
\begin{enumerate}[label=\alph*.]
\item HCl donates a proton to OH$^-$, forming water as Na$^+$ pairs with Cl$^-$
\item NaOH donates protons to HCl, forming hydrogen gas
\item Both reactants act as proton donors, preventing neutralization
\item Cl$^-$ acts as the proton donor while Na$^+$ accepts the proton
\end{enumerate}
\end{minipage}\\[0.2cm]
\hline
\end{tabular}

\noindent\begin{tabular}{|p{0.9\linewidth}|}
\hline
\textbf{(32)} In the reaction $\text{NH}_3 + \text{H}_2\text{O} \rightarrow \text{NH}_4^+ + \text{OH}^-$, which species acts as the Brönsted-Lowry base?\\[0.2cm]
\begin{minipage}[t]{0.85\linewidth}
\begin{enumerate}[label=\alph*.]
\item NH$_3$
\item H$_2$O
\item NH$_4^+$
\item OH$^-$
\end{enumerate}
\end{minipage}\\[0.2cm]
\hline
\end{tabular}
\newpage

\noindent\begin{tabular}{|p{0.9\linewidth}|}
\hline
\textbf{(33)} Identify the reaction type:\\[0.2cm]
2H$_2$O$_2$ $\rightarrow$ 2H$_2$O + O$_2$\\[0.2cm]
\begin{minipage}[t]{0.85\linewidth}
\begin{enumerate}[label=\alph*.]
\item Combination
\item Combustion
\item Decomposition
\item Exchange
\end{enumerate}
\end{minipage}\\[0.2cm]
\hline
\end{tabular}

\noindent\begin{tabular}{|p{0.9\linewidth}|}
\hline
\textbf{(34)} Which pair represents a conjugate acid-base pair?\\[0.2cm]
\begin{minipage}[t]{0.85\linewidth}
\begin{enumerate}[label=\alph*.]
\item NH$_3$/NH$_4^+$
\item OH$^-$/O$^{2-}$
\item Na$^+$/Na
\item Cl$^-$/Cl$_2$
\end{enumerate}
\end{minipage}\\[0.2cm]
\hline
\end{tabular}

\noindent\begin{tabular}{|p{0.9\linewidth}|}
\hline
\textbf{(35)} Which of the following elements always has an oxidation number of +1 in compounds?\\[0.2cm]
\begin{minipage}[t]{0.85\linewidth}
\begin{enumerate}[label=\alph*.]
\item Sodium
\item Oxygen
\item Fluorine
\item Sulfur
\end{enumerate}
\end{minipage}\\[0.2cm]
\hline
\end{tabular}

\noindent\begin{tabular}{|p{0.9\linewidth}|}
\hline
\textbf{(36)} What is the oxidation number of manganese in KMnO$_4$?\\[0.2cm]
\begin{minipage}[t]{0.85\linewidth}
\begin{enumerate}[label=\alph*.]
\item +2
\item +4
\item +6
\item +7
\end{enumerate}
\end{minipage}\\[0.2cm]
\hline
\end{tabular}
\newpage

\noindent\begin{tabular}{|p{0.9\linewidth}|}
\hline
\textbf{(37)} In a redox reaction, the reducing agent is the substance that:\\[0.2cm]
\begin{minipage}[t]{0.85\linewidth}
\begin{enumerate}[label=\alph*.]
\item Gains electrons
\item Loses electrons
\item Is neither oxidized nor reduced
\item Lowers its oxidation number
\end{enumerate}
\end{minipage}\\[0.2cm]
\hline
\end{tabular}

\noindent\begin{tabular}{|p{0.9\linewidth}|}
\hline
\textbf{(38)} For the reaction $2\text{SO}_2(g) + \text{O}_2(g) \rightleftharpoons 2\text{SO}_3(g)$, decreasing pressure favors: \\[0.2cm]
\begin{minipage}[t]{0.85\linewidth}
\begin{enumerate}[label=\alph*.]
\item Formation of SO$_3$
\item Formation of SO$_2$ and O$_2$
\item No change in equilibrium
\item Formation of whichever side absorbs heat
\end{enumerate}
\end{minipage}\\[0.2cm]
\hline
\end{tabular}


\noindent\begin{tabular}{|p{0.9\linewidth}|}
\hline
\textbf{(39)} For the endothermic reaction $\text{heat} + \text{H}_2 + \text{I}_2 \rightarrow 2\text{HI}$, decreasing temperature shifts equilibrium: \\[0.2cm]
\begin{minipage}[t]{0.85\linewidth}
\begin{enumerate}[label=\alph*.]
\item Toward products
\item Toward reactants
\item No shift
\item Depends only on pressure
\end{enumerate}
\end{minipage}\\[0.2cm]
\hline
\end{tabular}

\noindent\begin{tabular}{|p{0.9\linewidth}|}
\hline
\textbf{(40)} For the reaction $\text{A} + \text{B} \rightleftharpoons \text{C}$, if $K_{eq} = 12$, which statement is true?\\[0.2cm]
\begin{minipage}[t]{0.85\linewidth}
\begin{enumerate}[label=\alph*.]
\item Products are favored
\item Reactants are favored
\item Neither side is favored
\item The reaction goes to completion
\end{enumerate}
\end{minipage}\\[0.2cm]
\hline
\end{tabular}
\newpage

\noindent\begin{tabular}{|p{0.9\linewidth}|}
\hline
\textbf{(41)} What type of reaction occurs when lithium sulfate reacts with barium nitrate to form barium sulfate and lithium nitrate?\\[0.2cm]
Li$_2$SO$_4$ + Ba(NO$_3$)$_2$ $\rightarrow$ BaSO$_4$ + 2LiNO$_3$\\[0.2cm]
\begin{minipage}[t]{0.85\linewidth}
\begin{enumerate}[label=\alph*.]
\item Combination
\item Decomposition
\item Exchange
\item Combustion
\end{enumerate}
\end{minipage}\\[0.2cm]
\hline
\end{tabular}

\noindent\begin{tabular}{|p{0.9\linewidth}|}
\hline
\textbf{(42)} For the reaction $2\text{X} \rightleftharpoons \text{Y}$, $K_{eq} = 1.0$. Which is true at equilibrium?\\[0.2cm]
\begin{minipage}[t]{0.85\linewidth}
\begin{enumerate}[label=\alph*.]
\item $[Y]$ = [X]
\item $[Y]$ = $[X]^2$
\item $[X]$ = 2[Y]
\item $K_{eq}$ cannot equal 1 for reversible reactions
\end{enumerate}
\end{minipage}\\[0.2cm]
\hline
\end{tabular}

\noindent\begin{tabular}{|p{0.9\linewidth}|}
\hline
\textbf{(43)} What kind of reaction is the burning of glucose in oxygen to form carbon dioxide and water?\\[0.2cm]
C$_6$H$_{12}$O$_6$ + 6O$_2$ $\rightarrow$  6CO$_2$ + 6H$_2$O\\[0.2cm]
\begin{minipage}[t]{0.85\linewidth}
\begin{enumerate}[label=\alph*.]
\item Double displacement
\item Combustion
\item Decomposition
\item Combination
\end{enumerate}
\end{minipage}\\[0.2cm]
\hline
\end{tabular}

\noindent\begin{tabular}{|p{0.9\linewidth}|}
\hline
\textbf{(44)} A substance that gains electrons during a reaction is:\\[0.2cm]
\begin{minipage}[t]{0.85\linewidth}
\begin{enumerate}[label=\alph*.]
\item Oxidized
\item Reduced
\item Neutralized
\item Dissociated
\end{enumerate}
\end{minipage}\\[0.2cm]
\hline
\end{tabular}
\newpage

\noindent\begin{tabular}{|p{0.9\linewidth}|}
\hline
\textbf{(45)} Which of the following represents a redox reaction?\\[0.2cm]
\begin{minipage}[t]{0.85\linewidth}
\begin{enumerate}[label=\alph*.]
\item HCl + NaOH $\rightarrow$ NaCl + H$_2$O\\[0.1cm]
\item Zn + 2HCl $\rightarrow$ ZnCl$_2$ + H$_2$\\[0.1cm]
\item CaCO$_3$ $\rightarrow$ CaO + CO$_2$\\[0.1cm]
\item NaOH + HNO$_3$ $\rightarrow$ NaNO$_3$ + H$_2$O\\[0.1cm]
\end{enumerate}

\end{minipage}\\[0.2cm]
\hline
\end{tabular}

\noindent\begin{tabular}{|p{0.9\linewidth}|}
\hline
\textbf{(46)} For the exothermic decomposition $2\text{NO}_{(g)} \rightarrow {\text{N}_{2}}_{(g)} + {\text{O}_{2}}_{(g)} + \text{heat}$, adding heat will favor: \\[0.2cm]
\begin{minipage}[t]{0.85\linewidth}
\begin{enumerate}[label=\alph*.]
\item Formation of N$_2$ and O$_2$
\item Formation of NO
\item No change
\item Whichever side has fewer gas particles
\end{enumerate}
\end{minipage}\\[0.2cm]
\hline
\end{tabular}

\noindent\begin{tabular}{|p{0.9\linewidth}|}
\hline
\textbf{(47)} What type of reaction occurs when silver nitrate reacts with sodium chloride to produce silver chloride and sodium nitrate?\\[0.2cm]
AgNO$_3$ + NaCl $\rightarrow$ AgCl + NaNO$_3$\\[0.2cm]
\begin{minipage}[t]{0.85\linewidth}
\begin{enumerate}[label=\alph*.]
\item Combination
\item Decomposition
\item Exchange
\item Combustion
\end{enumerate}
\end{minipage}\\[0.2cm]
\hline
\end{tabular}

\noindent\begin{tabular}{|p{0.9\linewidth}|}
\hline
\textbf{(48)} For the reaction $\text{H}_2 + \text{I}_2 \rightleftharpoons 2\text{HI}$, which change increases $Q$?\\[0.2cm]
\begin{minipage}[t]{0.85\linewidth}
\begin{enumerate}[label=\alph*.]
\item Adding HI
\item Removing HI
\item Adding H$_2$
\item Removing heat
\end{enumerate}
\end{minipage}\\[0.2cm]
\hline
\end{tabular}
\newpage

\noindent\begin{tabular}{|p{0.9\linewidth}|}
\hline
\textbf{(49)} What type of chemical reaction does ethanol undergo in air?\\[0.2cm]
C$_2$H$_5$OH + 3O$_2$ $\rightarrow$ 2CO$_2$ + 3H$_2$O\\[0.2cm]
\begin{minipage}[t]{0.85\linewidth}
\begin{enumerate}[label=\alph*.]
\item Neutralization
\item Single replacement
\item Combustion
\item Precipitation
\end{enumerate}
\end{minipage}\\[0.2cm]
\hline
\end{tabular}

\noindent\begin{tabular}{|p{0.9\linewidth}|}
\hline
\textbf{(50)} For the reaction $\text{A} + \text{B} + \text{heat} \rightarrow \text{C}$, which way does equilibrium shift when heat is added?\\[0.2cm]
\begin{minipage}[t]{0.85\linewidth}
\begin{enumerate}[label=\alph*.]
\item Toward products
\item Toward reactants
\item No change
\item It depends on pressure only
\end{enumerate}
\end{minipage}\\[0.2cm]
\hline
\end{tabular}

\noindent\begin{tabular}{|p{0.9\linewidth}|}
\hline
\textbf{(51)} For the exothermic reaction $\text{X} + \text{Y} \rightarrow \text{Z} + \text{heat}$, removing heat shifts equilibrium: \\[0.2cm]
\begin{minipage}[t]{0.85\linewidth}
\begin{enumerate}[label=\alph*.]
\item Toward reactants
\item Toward products
\item Causes no shift
\item Only affects rate, not equilibrium
\end{enumerate}
\end{minipage}\\[0.2cm]
\hline
\end{tabular}

\noindent\begin{tabular}{|p{0.9\linewidth}|}
\hline
\textbf{(52)} In the equilibrium $\text{HF} + \text{H}_2\text{O} \rightleftharpoons \text{H}_3\text{O}^+ + \text{F}^-$, which is the conjugate acid of F$^-$? \\[0.2cm]
\begin{minipage}[t]{0.85\linewidth}
\begin{enumerate}[label=\alph*.]
\item HF
\item H$_3$O$^+$
\item H$_2$O
\item F$^-$ has no conjugate acid
\end{enumerate}
\end{minipage}\\[0.2cm]
\hline
\end{tabular}
\newpage

\noindent\begin{tabular}{|p{0.9\linewidth}|}
\hline
\textbf{(53)} What type of reaction occurs when acetylene reacts with oxygen to produce CO$_2$ and H$_2$O?\\[0.2cm]
2C$_2$H$_2$ + 5O$_2$ $\rightarrow$ 4CO$_2$ + 2H$_2$O\\[0.2cm]
\begin{minipage}[t]{0.85\linewidth}
\begin{enumerate}[label=\alph*.]
\item Neutralization
\item Combustion
\item Combination
\item Decomposition
\end{enumerate}
\end{minipage}\\[0.2cm]
\hline
\end{tabular}

\noindent\begin{tabular}{|p{0.9\linewidth}|}
\hline
\textbf{(54)} For the reaction $\text{N}_2(g) + 3\text{H}_2(g) \rightleftharpoons 2\text{NH}_3(g)$, increasing pressure shifts the equilibrium: \\[0.2cm]
\begin{minipage}[t]{0.85\linewidth}
\begin{enumerate}[label=\alph*.]
\item Toward products because fewer gas moles are present
\item Toward reactants because more gas moles are present
\item No shift because pressure does not affect gases
\item Toward whichever side has more volume
\end{enumerate}
\end{minipage}\\[0.2cm]
\hline
\end{tabular}

\noindent\begin{tabular}{|p{0.9\linewidth}|}
\hline
\textbf{(55)} In the reaction Zn + Cu$^{2+}$ $\rightarrow$ Zn$^{2+}$ + Cu, which species is oxidized?\\[0.2cm]
\begin{minipage}[t]{0.85\linewidth}
\begin{enumerate}[label=\alph*.]
\item Cu$^{2+}$
\item Cu
\item Zn
\item H$_2$O
\end{enumerate}
\end{minipage}\\[0.2cm]
\hline
\end{tabular}

\noindent\begin{tabular}{|p{0.9\linewidth}|}
\hline
\textbf{(56)} Identify the type of reaction below:\\[0.2cm]
K$_2$SO$_4$ + Ba(NO$_3$)$_2$ $\rightarrow$ 2KNO$_3$ + BaSO$_4$\\[0.2cm]
\begin{minipage}[t]{0.85\linewidth}
\begin{enumerate}[label=\alph*.]
\item Exchange
\item Combination
\item Combustion
\item Decomposition
\end{enumerate}
\end{minipage}\\[0.2cm]
\hline
\end{tabular}
\newpage

\noindent\begin{tabular}{|p{0.9\linewidth}|}
\hline
\textbf{(57)} For the reaction $\text{A} + \text{C} \rightleftharpoons \text{AC}$ with $K_{eq} = 0.02$, which is true?\\[0.2cm]
\begin{minipage}[t]{0.85\linewidth}
\begin{enumerate}[label=\alph*.]
\item Reactants are strongly favored
\item Products are strongly favored
\item Equal amounts of A and AC form
\item $K_{eq}$ must be recalculated using solids only
\end{enumerate}
\end{minipage}\\[0.2cm]
\hline
\end{tabular}

\noindent\begin{tabular}{|p{0.9\linewidth}|}
\hline
\textbf{(58)} Which of the following compounds acts as an Arrhenius base when dissolved in water?\\[0.2cm]
\begin{minipage}[t]{0.85\linewidth}
\begin{enumerate}[label=\alph*.]
\item NaOH, because it increases the concentration of OH$^-$ in solution
\item HCl, because it donates protons to other species
\item NH$_3$, because it donates an electron pair to H$^+$
\item CO$_2$, because it produces H$^+$ in water
\end{enumerate}
\end{minipage}\\[0.2cm]
\hline
\end{tabular}


\noindent\begin{tabular}{|p{0.9\linewidth}|}
\hline
\textbf{(59)} For $\text{D} \rightleftharpoons \text{E}$, which scenario decreases the value of the reaction quotient $Q$?\\[0.2cm]
\begin{minipage}[t]{0.85\linewidth}
\begin{enumerate}[label=\alph*.]
\item Adding D
\item Removing D
\item Adding E
\item Both b. and c.
\end{enumerate}
\end{minipage}\\[0.2cm]
\hline
\end{tabular}

\noindent\begin{tabular}{|p{0.9\linewidth}|}
\hline
\textbf{(60)} For the reaction $\text{M} \rightleftharpoons \text{N}$, if $K_{eq}$ is extremely large, which is true?\\[0.2cm]
\begin{minipage}[t]{0.85\linewidth}
\begin{enumerate}[label=\alph*.]
\item Products dominate at equilibrium
\item Reactants dominate at equilibrium
\item $[M]$ must be greater than [N]
\item Q must equal zero
\end{enumerate}
\end{minipage}\\[0.2cm]
\hline
\end{tabular}
\end{document}



